\section{OLLM}
\label{sec:ollm}

\ldeen{OLLM ist die ausführbare Buildschnittstelle des Bundles. Es liest das
Manifest, entdeckt Units, normalisiert Target und Sprache, erzeugt einen
isolierten Buildauftrag und lässt ihn von \texttt{latexmk} mit Lua\LaTeX\
ausführen. Dadurch müssen Quelldateien weder Jobnamen zerlegen noch den
Projektzustand erraten. Die allgemeine Befehlsform lautet}{OLLM is the
executable build interface of the bundle. It reads the manifest, discovers
units, normalizes the target and language, creates an isolated build request,
and runs it through \texttt{latexmk} with Lua\LaTeX. Source files therefore do
not need to parse job names or infer project state. The general command form is}:

\begin{verbatim}
ollm [globale Optionen] [Aktion] [Target] [Aktionsoptionen]
\end{verbatim}

\ldeen{Ohne Aktion wird \texttt{build} angenommen. Die mitgelieferten Targets
\texttt{slides}, \texttt{handout} und \texttt{script} dürfen als nacktes Wort
stehen; projektdefinierte Targets werden eindeutig mit
\texttt{--target=NAME} gewählt}{If no action is given, \texttt{build} is
assumed. The supplied targets \texttt{slides}, \texttt{handout}, and
\texttt{script} may be written as bare words; project-defined targets are
selected unambiguously with \texttt{--target=NAME}}:

\begin{verbatim}
ollm slides --language=de
ollm build --target=script --language=en
\end{verbatim}

\subsection{\ldeen{Bauen: von einer Instanz bis zur Serie}{Building: From One Instance to a Series}}

\ldeen{Für die tägliche Arbeit in einer Unit genügt meist ein einzelner
Build. Ohne \texttt{--language} verwendet OLLM die Projektsprache, sofern das
Target sie unterstützt, sonst dessen erste konfigurierte Sprache.
\texttt{--source} ist standardmäßig \texttt{main.tex}}{For day-to-day work in
a unit, a single build is usually sufficient. Without \texttt{--language},
OLLM uses the project default language if supported by the target, otherwise
the target's first configured language. \texttt{--source} defaults to
\texttt{main.tex}}:

\begin{verbatim}
ollm build script --language=en
ollm build --target=slides --source=main.tex
\end{verbatim}

\ldeen{Der Scope bestimmt, wie weit die Auswahl reicht. \texttt{current} baut
eine Projektion der aktuellen Unit. \texttt{unit} baut deren zulässige
Target-/Sprachmatrix. \texttt{series} baut eine Projektion aller passenden
normalen Units, während \texttt{collection} die eindeutig passende
Integrationsunit des Dokumenttyps baut. Die letzten beiden Scopes funktionieren
auch aus der Projektwurzel}{The scope determines how far the selection reaches.
\texttt{current} builds one projection of the current unit. \texttt{unit}
builds its permitted target/language matrix. \texttt{series} builds one
projection of every applicable regular unit, whereas \texttt{collection}
builds the uniquely matching integration unit for the document type. The last
two scopes also work from the project root}:

\begin{verbatim}
ollm build --all --scope=unit
ollm handout --language=de --scope=series
ollm script --language=de --scope=collection
\end{verbatim}

\ldeen{\texttt{--all} erweitert die Matrix nur innerhalb des gewählten
Scopes; es bedeutet nicht automatisch "`das ganze Projekt"'.
\texttt{--rebuild} erzwingt einen vollständigen \texttt{latexmk}-Neubau.
\texttt{--dry-run} zeigt die normalisierte Auswahl, ohne \LaTeX\ zu starten,
und eignet sich zusammen mit \texttt{--format=json} für Werkzeuge und CI.}{
\texttt{--all} expands the matrix only within the selected scope; it does not
automatically mean "`the entire project"'. \texttt{--rebuild} forces a full
\texttt{latexmk} rebuild. \texttt{--dry-run} displays the normalized selection
without starting \LaTeX\ and, together with \texttt{--format=json}, is useful
for tools and CI.}

\ldeen{Serieninterne Referenzen und Integrationen bilden Abhängigkeiten.
\texttt{--resolve} baut veraltete Produzenten und Konsumenten in mehreren
Runden bis zu einem Fixpunkt. Eine logische Unit, die noch nie direkt gebaut
wurde, muss zunächst einmal in der benötigten Projektion gebaut werden, damit
OLLM ihre ID dem Verzeichnis zuordnen kann.}{Cross-unit references and
integrations form dependencies. \texttt{--resolve} rebuilds stale producers
and consumers in several rounds until a fixed point is reached. A logical unit
that has never been built directly must first be built once in the required
projection so that OLLM can associate its ID with its directory.}

\begin{verbatim}
ollm build script --language=de --resolve
\end{verbatim}

\subsection{\ldeen{Prüfen und Diagnostizieren}{Inspecting and Diagnosing}}

\ldeen{\texttt{report}, \texttt{check} und \texttt{doctor} beantworten drei
verschiedene Fragen. \texttt{report} beschreibt den bekannten Zustand:
Units, Projektionen, Generationen, Artefakte und Abhängigkeiten. Befunde ändern
seinen Exitstatus nicht, solange der Zustand lesbar ist. \texttt{check} liest
dieselben Daten, bewertet Inkonsistenzen aber als Fehler und eignet sich daher
als Freigabegate. Beide Befehle bauen nichts.}{\texttt{report},
\texttt{check}, and \texttt{doctor} answer three different questions.
\texttt{report} describes the known state: units, projections, generations,
artifacts, and dependencies. Findings do not change its exit status as long as
the state is readable. \texttt{check} reads the same data but treats
inconsistencies as errors and is therefore suitable as a release gate. Neither
command builds anything.}

\begin{verbatim}
ollm report --scope=series
ollm check --scope=current --target=script --language=de
ollm check --format=json
\end{verbatim}

\ldeen{\texttt{doctor} prüft dagegen die Werkzeugkette: Perl,
\texttt{latexmk}, Lua\LaTeX, \texttt{kpsewhich}, den TOML-Parser und die
benötigten TeX-Pakete. In einem Projekt kommen Manifest, Definitionen,
Schreibbarkeit, Zustand und Shell-Escape-Policy hinzu. Verwenden Sie
\texttt{doctor}, wenn bereits die Umgebung fraglich ist, und \texttt{check},
wenn ein vorhandener Projektzustand widersprüchlich erscheint.}{
\texttt{doctor}, by contrast, checks the toolchain: Perl, \texttt{latexmk},
Lua\LaTeX, \texttt{kpsewhich}, the TOML parser, and required TeX packages.
Inside a project it also checks the manifest, definitions, writability, state,
and shell-escape policy. Use \texttt{doctor} when the environment itself is in
doubt, and \texttt{check} when an existing project state appears inconsistent.}

\begin{verbatim}
ollm doctor
ollm doctor --project-root=/path/to/project --format=json
\end{verbatim}

\subsection{\ldeen{Aufräumen und Zustände pflegen}{Cleaning and Maintaining State}}

\ldeen{OLLM trennt Buildverzeichnisse von promoviertem Serienzustand.
\texttt{clean --level=aux} entfernt Hilfsdateien, behält aber das lokale PDF;
\texttt{build} entfernt das isolierte Buildverzeichnis; \texttt{state}
entfernt promovierte Projektionen; \texttt{all} kombiniert die letzten beiden.
Mit \texttt{--dry-run} sehen Sie zuerst den Löschplan. Scopes sind hier
\texttt{current}, \texttt{unit} und \texttt{series}.}{OLLM separates build
directories from promoted series state. \texttt{clean --level=aux} removes
auxiliary files but retains the local PDF; \texttt{build} removes the isolated
build directory; \texttt{state} removes promoted projections; and
\texttt{all} combines the last two. Use \texttt{--dry-run} to inspect the
deletion plan first. The available scopes here are \texttt{current},
\texttt{unit}, and \texttt{series}.}

\begin{verbatim}
ollm clean --level=aux
ollm clean --level=build --scope=unit --dry-run
ollm clean --level=state --scope=series
\end{verbatim}

\ldeen{\texttt{prune} beseitigt aufgegebene temporäre Zustände und ersetzte
Generationen. Zustände physisch entfernter Units werden nur mit
\texttt{--stale-units} einbezogen. Das ist absichtlich vorsichtiger als ein
pauschales Löschen von \texttt{.osglecture}, bei dem auch die bekannten
Zuordnungen logischer Unit-IDs verloren gehen.}{\texttt{prune} removes
abandoned temporary states and superseded generations. State belonging to
physically removed units is included only with \texttt{--stale-units}. This is
deliberately more cautious than deleting \texttt{.osglecture} wholesale,
which also loses known mappings of logical unit IDs.}

\begin{verbatim}
ollm prune --dry-run
ollm prune --stale-units
\end{verbatim}

\subsection{\ldeen{Bereitstellen}{Deployment}}

\ldeen{\texttt{deploy} kopiert bereits erfolgreich promovierte PDFs an die im
Manifest konfigurierten Zielorte. Es startet keinen Build und legt keine
Zielverzeichnisse an. Verwenden Sie daher zuerst \texttt{build} oder
\texttt{build --resolve} und danach \texttt{deploy}. In der Projektwurzel gilt
standardmäßig Scope \texttt{series}, in einer Unit \texttt{current}.}{
\texttt{deploy} copies already successfully promoted PDFs to destinations
configured in the manifest. It neither starts a build nor creates destination
directories. Therefore run \texttt{build} or \texttt{build --resolve} first,
then \texttt{deploy}. The default scope is \texttt{series} at the project root
and \texttt{current} inside a unit.}

\begin{verbatim}
ollm deploy handout --language=de --scope=series
ollm deploy --target=script --scope=collection
ollm deploy --all --overwrite
\end{verbatim}

\ldeen{Dateinamenschablonen können unter anderem Serien-ID, Unit-ID,
Ordnungskennung, Kapitel, Dokumenttyp, Sprache und Rolle verwenden. Eine
vorhandene identische Datei bleibt unverändert. Ob eine abweichende Datei
\texttt{--overwrite} verlangt, bestimmt
\texttt{security.deployment.overwrite}.}{File-name templates can use, among
other values, the series ID, unit ID, ordinal, chapter, document type,
language, and role. An existing identical file remains untouched. Whether a
different existing file requires \texttt{--overwrite} is controlled by
\texttt{security.deployment.overwrite}.}

\subsection{\ldeen{Projekt anlegen und allgemeine Optionen}{Creating a Project and General Options}}

\ldeen{\texttt{ollm newtoml} erzeugt ein neues
\texttt{ollmconfig.toml}, überschreibt aber kein bestehendes Manifest. Für die
genaue Syntax einer installierten Version sind \texttt{ollm --help} und
\texttt{ollm --version} maßgeblich. Häufige globale Optionen sind
\texttt{--config=DATEI}, \texttt{--project-root=VERZEICHNIS},
\texttt{--format=text|json}, \texttt{--color=auto|always|never} und
\texttt{--non-interactive}.}{\texttt{ollm newtoml} creates a new
\texttt{ollmconfig.toml} but never overwrites an existing manifest. For the
exact syntax of an installed version, \texttt{ollm --help} and
\texttt{ollm --version} are authoritative. Common global options are
\texttt{--config=FILE}, \texttt{--project-root=DIRECTORY},
\texttt{--format=text|json}, \texttt{--color=auto|always|never}, and
\texttt{--non-interactive}.}

\ldeen{Für Builddiagnosen wählen \texttt{--debug=tex}, \texttt{ollm} oder
\texttt{tex+ollm} die Ausgabeschicht; \texttt{--warnings=all},
\texttt{important} oder \texttt{none} steuert die Warnungsmenge. Unbekannte
kompatible Minusoptionen werden beim Build an \texttt{latexmk} weitergereicht;
Optionen, die den kontrollierten Buildvertrag umgehen würden, weist OLLM
zurück.}{For build diagnostics, \texttt{--debug=tex}, \texttt{ollm}, or
\texttt{tex+ollm} selects the output layer; \texttt{--warnings=all},
\texttt{important}, or \texttt{none} controls warning verbosity. Compatible
unknown dash options are forwarded to \texttt{latexmk} during a build; OLLM
rejects options that would bypass the controlled build contract.}
